\section{Discussion}\label{sec:discussion}

PLACEHOLDER: Two practical conclusions emerge for operating a solar portfolio under forecast uncertainty.

First, \emph{the temporal structure of error is partly correctable.} The dominant pattern in
error growth is a diurnal cycle (statistically significant for GHI, RH, U10, and wind direction),
not monotonic horizon-driven drift. A harmonic bias correction keyed to time of day would remove
a meaningful component of error without any change to the underlying forecast model. The fixed
\texttt{06z} initialization limits how cleanly we can separate genuine lead-time growth from the
diurnal cycle; forecasts from multiple initialization times would be required to fully resolve
that confound.

Second, \emph{spatial diversification is real but variable-specific and season-specific.}
Cloud-driven irradiance error---the component that matters most for solar output---decorrelates
over the shortest distances, so a geographically spread fleet genuinely smooths it, especially in
summer. Synoptic fields such as surface pressure stay correlated across the entire PJM footprint
and cannot be diversified away. Any portfolio-level uncertainty model should therefore treat these
variable classes differently rather than assuming a single fleet-wide correlation.
